\scanpage{12}
\uline{Exercice} \handwritten{4}.\handwritten{10}.\handwritten{7} . On reprend les hypothèses et les notations de l'exercice \handwritten{4}.\handwritten{10}.\handwritten{6} . On suppose de plus que \(\underline{M}\) possède la propriété suivante :

(Et) Pour tout couple de flèches composables \(u\) et \(v\) de \(\mathrm{Fl}\ S\) , les \deleted{ces} propriétés \(u\circ v\unsure{\varepsilon|\in}\underset{\handwritten{\deleted{\unsure{m|M}}}}{\underline{M}}\) et \(u\varepsilon\underline{M}\) impliquent \(v\handwritten{\unsure{\in|\varepsilon}}\underline{M}\) .

1) Montrer que pour toute flèche \(u:X\to Y\) , le foncteur de changement de base par \(u\) est un morphisme \deleted{du topos} \(S(u)\) du site \(S(X)\) dans le site \(S(Y)\) . En d'autres termes , montrer que \(S(u)^{\handwritten{*}}:S(Y)\to S(X)\) est exact à gauche .

2) Montrer que \(\operatorname{Prol}_X:S(X)^{\handwritten{\sim}}\to(S_{\handwritten{/}X})^{\handwritten{\sim}}\) est exact à gauche . En déduire qu'il existe un morphisme de topos
\[
g:S^{\handwritten{\sim}}_{\handwritten{\deleted{\unsure{u|m}}/}X}\to S(X)^{\handwritten{\sim}}
\]
tel que \(g^{\handwritten{*}}=\operatorname{Prol}_X\) et \(g_{\handwritten{*}}=\operatorname{Res}_X\) . Montrer que \(g\) est une rétraction du plongement \(f\) i.e qu'on a \(g\circ f\handwritten{\simeq}\operatorname{id}_{S(X)}\) et qu'on a un morphisme ``d'adjonction'' (entre morphismes de topos ) \(\operatorname{id}_{S^{\handwritten{\sim}}_{\handwritten{/}X}}\handwritten{\longrightarrow} f\circ g\) .

\handwritten{3}) Montrer que les conclusion de \handwritten{1}) et \handwritten{2}) ne sont plus valables lorsqu'on ne fait pas l'hypothèse (Et) . (\deleted{p}Prendre pour \(S\) la catégorie des schémas \unsure{munie|\deleted{,}munie} de la topologie f.p.p.f. (fidèlement plate et de présentation finie \handwritten{[} S.G.A. \handwritten{3} IV \handwritten{]} ) et pour \(\underline{M}\) les morphismes \deleted{fidèlement} plats et de \unsure{p\'{r}ésentation|présentation} finie ).

\handwritten{4}) Acheter une médaille en chocolat pour les rédacteurs .
% archive sheet number: 11
